\section{Introduction}\label{sec:introduction}

Near a reducible polystable bundle, deformation theory gives a finite
quotient description of the moduli germ. The harmonic $L^2$ metric
contains further information. It is obtained by solving the
Hermitian--Einstein equation and then taking harmonic representatives
of deformations. When the stabilizer changes, the finite orbit
operator develops small eigenvalues, so an inverse estimate that
treats all gauge directions together is not uniform. This distinction
is especially relevant when several stability parameters approach
different walls at different rates.

We prove a local comparison between the actual harmonic metric and
the flat quiver quotient metric that remains relative through these
degenerations. The geometric setting is a projective K3 surface with
a polystable locally free sheaf
$E_0=\bigoplus_iF_i\otimes\CC^{m_i}$.
The factors are pairwise nonisomorphic slope-stable bundles with
equal reduced Hilbert polynomial for an integral ample class.
The central analytic metric is Ricci-flat in a K\"ahler class for
which the factors are stable of a common slope; the two classes
may coincide. These hypotheses are verified directly in the
application below, including at its selected noncentral wall points.

The central harmonic norm and Serre duality identify
$V=\Ext^1(E_0,E_0)$ with a unitary doubled-quiver representation.
Its complex moment zero locus gives the convergent integrability
model. For nearby K\"ahler forms $\omega_t$, the real central
parameter is the degree difference
\[
 \zeta_i(t)=2\pi\bigl(\rank(F_i)\mu_{\omega_t}(E_0)
              -\deg_{\omega_t}F_i\bigr).
\]
Theorem~\ref{thm:harmonic-comparison} identifies full local
neighborhoods of actual HYM solutions with the corresponding
finite quotients. On each genuine smooth-stratum tangent space,
at a representative of harmonic norm $r$, it proves
\[
 f_t^*g_t^{L^2}=(1+O(|t|+r))g_{\zeta(t)}
\]
in the sense of two-sided inequalities of Hermitian forms.
The constant does not depend on stabilizer type or ratios among
positive finite orbit eigenvalues. A verified central sign symmetry
improves $r$ to $r^2$; no such cancellation is assumed otherwise.

\subsection{The analytic argument}
All parallel automorphisms are retained in the finite-dimensional
problem. On their fixed infinite-dimensional orthogonal complement,
the gauge equation has a uniformly coercive elliptic inverse.
The remaining finite moment map is the moment of a smooth closed
invariant two-form close to the flat form.

The decisive finite estimate is relative to each orthogonal
vertex projection $P$:
\[
 |\langle\rho(x)-\mu(x),P\rangle|
       \le\varepsilon\|C_xP\|^2.
\]
Here $\mu$ is the flat moment, $\rho$ the perturbed normalized
moment, and $C_x$ the infinitesimal action. The same crossing
energy appears in the positive moment increment across a spectral
cut. Working in the Hermitian centralizer of the actual compact
stabilizer, orthogonal to its infinitesimal stabilizer, makes this
crossing energy strictly positive at every nonzero boundary
correction. It cancels from the estimate. A degree argument gives
a bounded complex-orbit correction and a reverse correction
without assuming a uniform positive orbit gap.

For the metric, first minimize over the fixed complementary
gauge directions. The resulting smooth finite Hermitian form
has the flat harmonic norm as its value at the origin. Minimizing
next over the retained finite algebra gives exactly the actual
harmonic Kodaira--Spencer norm. Derivatives of the orbit correction
are complex vertical and disappear in this infimum. This yields
relative quotient estimates without bounding those derivatives.
A separate stratum implicit-function argument and a full local
inverse supply the differentiability and coverage needed in this
calculation.

These steps develop methods from the preceding stability-degeneration
papers. The spectral-cut calibration is from
\cite{Wue26a}, and its repeated-factor form is developed in
\cite{Wue26c}. Relative moduli-metric comparison by complementary
minimization is developed in \cite{Wue26b,Wue26c}.
The projection-relative moment estimate and the harmonic
coordinate argument were used for the two-factor K3 example in
\cite{Wue26d}. Here the finite correction treats arbitrary
quiver stabilizers, including loops, and its reverse direction
covers the full nearby solution space. We give the required
proofs rather than import conclusions restricted to those earlier
representations.

\subsection{Intrinsic consequences}
A comparison of tangent norms alone does not determine intrinsic
distances or a complete rescaled limit: short competing paths
could leave the chart, and an image of a fixed model ball need
not cover the corresponding actual ball. We establish both
path localization and inverse coverage.

The finite-dimensional result is also useful independently.
For any doubled quiver whose complex zero fiber contains a
module-simple representation, meaning one with no nonzero proper
subrepresentation, Theorem~\ref{thm:varying-parameters}
proves that the full stratified length quotient at every real
central parameter is a proper geodesic space.
Arbitrary converging real parameters give converging distances
for converging balanced representatives, and convergence to the
whole complete pointed quotient, including singular levels.
The proof uses connected fibers of projective face maps to
construct shrinking neighborhoods, uniform semialgebraic paths
for short joins, and a quantitative Hamiltonian slice to retain
asymptotically exact lengths along strata. Horizontal lifting
gives the opposite distance inequality.

Combining this result with the harmonic comparison gives complete
actual rescaled limits whenever $\zeta(t_n)/\delta_n^2$ converges.
At a moving wall point, the full normal Ext-quiver again has a
module-simple zero representation. Translation in its trivial
factor and another rescaling identify finer limits.
The result can be iterated through any fixed finite succession
of convergent centers and scales, under the explicit bounded
normal-parameter and base-point hypotheses of
Theorem~\ref{thm:moving-centers}. No comparability of the successive
positive scales is imposed.
We also obtain relative distances to specified conical strata,
the square-root contraction profile of compact cores, and
matching to the surrounding free wall metric in the regime
$\|\zeta\|/\lambda\to0$, where $\lambda$ is the least positive
finite orbit eigenvalue.

\subsection{The K3 application}
We realize a projective K3 surface with
$\operatorname{NS}(X)=\langle2\rangle\oplus A_2(-2)$ and an ample
class $H$. Three spherical line bundles with classes $a,b,-a-b$
have four cross-extension directions for every pair. With
multiplicities $(2,2,1)$ their direct sum has primitive Mukai
vector $(5,a+b,-5)$ of square $46$.
The resulting local moduli have complex dimension $48$.
The real polarization family $H+ta+ub$ realizes two independent
stability parameters.

The finite representation is a doubled triangle on which two
matrix groups $U(2)$ act nontrivially and interact.
Five realized wall lines meet at the central bundle and divide
the parameter plane into ten chambers.
Explicit adjoint-paired arrow matrices prove the existence of simple
representations at zero moment level for the full vector and every required wall
summand. At selected actual two-factor wall strata, the transverse
contractions have models $T^*\PP^9$, $T^*\PP^{11}$ and
$T^*\PP^{17}$, with complete product metric limits and precise
square-root contraction scales. At the wall the distance to the
two-factor stratum is $(1+O(r))\sqrt{2s}$, where $s$ is the nonzero
singular value of the rank-one square-zero transverse invariant.
Theorem~\ref{thm:k3-application} states these conclusions for the
actual harmonic metric.

The cycle, interacting matrix symmetries, and independent
parameters are essential features of this application.
There is no vertex-scalar sign acting as $-I$ on every edge,
and approaches to its intersecting walls require the singular
and iterated limit arguments. The comparison is therefore used
beyond a change of matrix size in a two-factor quotient.

\subsection{Relation to the literature and scope}
The quadratic K3 deformation model is classical in origin.
Budur--Zhang prove the required formality \cite{BZ19};
Arbarello--Sacc\`a give the equivariant analytic quadratic
identification and its quiver description
\cite[Propositions~4.4 and 6.1]{AS18}.
Their dimension-one sheaf wall-crossing theorem is not used
to identify this rank-five bundle problem.
Their later treatment of Bridgeland moduli compares nearby stability
conditions with local quiver VGIT, away from totally semistable walls
and for sufficiently close conditions with the same heart
\cite[\S6, Claim~6.1, Proposition~6.3 and Corollary~6.4]{AS25}.
This broader local comparison concerns algebraic stability and does
not supply the harmonic metric estimates studied here.
King's quotient theory \cite{King94}, Crawley-Boevey's
moment-fiber theorem \cite{CB01}, and Bellamy--Schedler's
normality and projective face maps \cite{BeSc24} supply the
algebraic geometry used here. Mayrand's full Hamiltonian slice
\cite{May22} supplies the analytic local form; we additionally
prescribe its isometric normal first jet and compare its
real quotient metric.

Buchdahl--Schumacher's analytic GIT and Weil--Petersson work
\cite{BS21,BS23,BS26} provides analytic metric antecedents.
Clarke--Tipler prove local HYM wall variation for simple sufficiently
smooth semistable bundles with abelian graded automorphism group
\cite[Proposition~1.1 and Theorem~1.2]{CT24}.
Delloque removes the abelian restriction and proves continuity of
HYM connection classes under metric variation
\cite[Theorems~4.8--4.9]{DelCont25}.
These results concern connections and graded bundles; the relative
harmonic moduli-metric estimates require the additional arguments here.
We use Delloque's later parameter-dependent local solution theory
\cite{Del25} on its integrable HYM branch for full local
classification and inverse coverage, and prove separately the
uniform estimates through the retained symmetry directions.
Namikawa--Odaka establish metric tangent-cone asymptotics under
their reduction hypotheses \cite[Setup~1, Proposition~6.10 and
Corollary~6.13]{NO26}. Their Remark~6.14 distinguishes
the metric tangent cone from the flat symplectic-slice reduction
metric, and the derivative estimate in Corollary~6.13(ii)
has an isolated-singularity hypothesis. Our conclusions concern
the stated quiver and bundle settings, with nonisolated strata
and varying real parameters; they are not a general hyperk\"ahler
equivalence theorem for arbitrary symplectic singularities.
Xu--Zhuang, with an appendix by Guenancia, prove formal conicality
and derive analytic conicality in their symplectic-singularity
setting \cite[Theorem~1.2 and Appendix~A]{XZ26}; general
hyperk\"ahler-cone existence remains their Conjecture~1.3.
These algebraic and analytic conclusions do not by themselves
give the quantitative harmonic-metric comparison proved here.

All moduli identifications and contraction maps asserted here
are local. In particular we do not identify a global metric
completion, a bubbling compactification, or an HYM evolution.
The full stratified length metric is specified explicitly;
identifying it with the completion of the regular locus alone
would be a further statement.

